\answer{ex:05:1}{$\ell\approx17\um$; $\approx100$~суток: рассылку по мозгу делает ветвление провода, диффузия лишь размывает каждое место выброса на $\sim20\um$.}
\answer{ex:05:2}{$r=ab/(1+G)$, $S=1/(1+G)$ точно; $+1.7\,\%$ вместо $+20\,\%$.}
\answer{ex:05:3}{$T^*=1.21\,\text{с}$ при $G=2$ и $0.19\,\text{с}$ при $G=9$: реально $G\sim2$--$4$, подавление в $3$--$5$ раз.}
\answer{ex:05:4}{$\mathrm{CV}=1/\sqrt{2\lambda\tau_c}\approx9\cdot10^{-4}$ при суммарной частоте $\lambda$; одному нейрону нужно $\tau_c=12.5\,\text{с}$.}
\answer{ex:05:5}{$c\approx0.40$ в обоих случаях; $\mathrm{CV}=0.050$ против $0.158$, в $\sqrt{10}$ раз. Частоты отдельных нейронов остаются пропорциональны $a_i$: регулируется только сумма.}
\answer{ex:05:6}{$N_1=\overline{A\vee B}$, $N_2=\overline{A\vee N_1}$, $N_3=\overline{B\vee N_1}$, $N_4=\overline{N_2\vee N_3}$, XOR $=N_5=\overline{N_4}$; каждое переключение занимает $\tau_c\ln2$, путь из $4$ вентилей --- $2.8\,\text{с}$ на один XOR.}
\answer{ex:05:7}{(а)~$\tau_\gamma=6/\ln3\approx5.5\,\text{с}$. (б)~$\bar D<4\,\text{с}$ против $D<2.2\,\text{с}$: непредсказуемость задержки делает терпеливее.}
\answer{ex:05:8}{$k_2=1/\tau_d=10\,\text{с}^{-1}$, $k_1=1/\tau_d-1/\tau_\gamma=9.5\,\text{с}^{-1}$; $\tau_\gamma=1/(k_2-mk_1)$: удвоение при $+2.6\,\%$ (а при $+5.3\,\%$ горизонт бесконечен).}
